% GENERATED FILE. Edit canonical lessons, then run npm run generate:books.
% canonical-source-hash: fnv1a64:2811213ab7df36ec
% canonical-lessons: MW-C181-jamavno, MW-C181-bilovno, MW-C181-mathani, MW-C181-galno, MW-C181-khir

\chapter{Curd and Butter}
\label{ch:mw-curd-and-butter}
\hlchaptermodality{181}

\begin{chapteropening}
\textbf{By the end of this chapter:} \emph{I can say how milk is set, churned and strained.}

Complete the last lesson of chapter 181: I can say how milk is set, churned and strained.
\end{chapteropening}

\section[\texorpdfstring{\mw{जमावणो} (jamāvṇo)}{जमावणो (jamāvṇo)}]{\mw{जमावणो} (jamāvṇo) — to set curd}

\label{lesson:MW-C181-jamavno}



\begin{quote}
Before the new one: say the Marwadi for a hanging net, then the Marwadi for a flat basket.
\end{quote}

\subsection*{You'll want to know: \mw{जमावणो}}
\textbf{\mw{जमावणो}} — \emph{jamāvṇo} — \textquotedblleft{}to set curd\textquotedblright{}.

\begin{quote}
\mw{मां} \mw{रात} \mw{मांय} \mw{दही} \mw{जमावै} \mw{है।}
\end{quote}

\begin{quote}
\emph{mā̃ rāt mā̃y dahī jamāvai hai.}
\end{quote}

\begin{quote}
Mother sets the curd at night.
\end{quote}

Past (I …): \mw{म्हैं} \mw{दही} \mw{जमायो}, \emph{mhaĩ dahī jamāyo}.

\begin{cousinweb}
\textbf{\mw{जमावणो}} (\emph{jamāvṇo}) is built on Sanskrit \emph{yamyate}, \textquotedblleft{}is held fast\textquotedblright{}; CDIAL lists Old Marwari \emph{jamaï}, \textquotedblleft{}is collected\textquotedblright{}, and Hindi \emph{jamānā} (CDIAL 10428). Milk that sets holds fast as curd.
\end{cousinweb}

\begin{grammarlens}[title={What you've built}]
One more word for this chapter.
\end{grammarlens}

\subsection*{Your turn}
\begin{itemize}
\raggedright
  \item \emph{Say it:} \emph{jamāvṇo}
  \item \emph{Say it:} \emph{jamāvṇo}, once more
  \item \emph{Say it:} \emph{mā̃ rāt mā̃y dahī jamāvai hai.} — Mother sets the curd at night.
  \item \emph{Recall:} say the Marwadi for a hanging net, then the Marwadi for a flat basket, then say \emph{jamāvṇo} again
\end{itemize}

\begin{tcolorbox}[breakable,colback=teal!4,colframe=teal!35!black,title={Before you move on}]
What does \mw{जमावणो} mean? (\textquotedblleft{}To set curd\textquotedblright{}.) Say it once more.
\end{tcolorbox}
\section[\texorpdfstring{\mw{बिलोवणो} (bilovṇo)}{बिलोवणो (bilovṇo)}]{\mw{बिलोवणो} (bilovṇo) — to churn}

\label{lesson:MW-C181-bilovno}



\begin{quote}
Before the new one: say the Marwadi for a flat basket, then the Marwadi for to set curd.
\end{quote}

\subsection*{You'll want to know: \mw{बिलोवणो}}
\textbf{\mw{बिलोवणो}} — \emph{bilovṇo} — \textquotedblleft{}to churn\textquotedblright{}.

\begin{quote}
\mw{दादी} \mw{दिनूगै} \mw{दही} \mw{बिलोवै} \mw{है।}
\end{quote}

\begin{quote}
\emph{dādī dinūgai dahī bilovai hai.}
\end{quote}

\begin{quote}
Grandmother churns the curd early in the morning.
\end{quote}

Past (I …): \mw{म्हैं} \mw{दही} \mw{बिलोयो}, \emph{mhaĩ dahī biloyo}.

\begin{cousinweb}
\textbf{\mw{बिलोवणो}} (\emph{bilovṇo}) is inherited from Sanskrit \emph{vilobhayati}, \textquotedblleft{}perplexes, stirs up\textquotedblright{}; CDIAL lists Old Marwari \emph{bilovaṇā}, \textquotedblleft{}churnings\textquotedblright{} (CDIAL 11913a).
\end{cousinweb}

\begin{grammarlens}[title={What you've built}]
One more word for this chapter.
\end{grammarlens}

\subsection*{Your turn}
\begin{itemize}
\raggedright
  \item \emph{Say it:} \emph{bilovṇo}
  \item \emph{Say it:} \emph{bilovṇo}, once more
  \item \emph{Say it:} \emph{dādī dinūgai dahī bilovai hai.} — Grandmother churns the curd early in the morning.
  \item \emph{Recall:} say the Marwadi for a flat basket, then the Marwadi for to set curd, then say \emph{bilovṇo} again
\end{itemize}

\begin{tcolorbox}[breakable,colback=teal!4,colframe=teal!35!black,title={Before you move on}]
What does \mw{बिलोवणो} mean? (\textquotedblleft{}To churn\textquotedblright{}.) Say it once more.
\end{tcolorbox}
\section[\texorpdfstring{\mw{मथाणी} (mathāṇī)}{मथाणी (mathāṇī)}]{\mw{मथाणी} (mathāṇī) — a churning stick}

\label{lesson:MW-C181-mathani}



\begin{quote}
Before the new one: say the Marwadi for to set curd, then the Marwadi for to churn.
\end{quote}

\subsection*{You'll want to know: \mw{मथाणी}}
\textbf{\mw{मथाणी}} — \emph{mathāṇī} — \textquotedblleft{}a churning stick\textquotedblright{}.

\begin{quote}
\mw{मथाणी} \mw{दही} \mw{मांय} \mw{है।}
\end{quote}

\begin{quote}
\emph{mathāṇī dahī mā̃y hai.}
\end{quote}

\begin{quote}
The churning stick is in the curd.
\end{quote}

A feminine noun.

\begin{cousinweb}
\textbf{\mw{मथाणी}} (\emph{mathāṇī}) is inherited from Sanskrit \emph{manthāna}, \textquotedblleft{}a churning stick\textquotedblright{}; CDIAL lists Old Marwari \emph{mathāṇī} (CDIAL 9842).
\end{cousinweb}

\begin{grammarlens}[title={What you've built}]
One more word for this chapter.
\end{grammarlens}

\subsection*{Your turn}
\begin{itemize}
\raggedright
  \item \emph{Say it:} \emph{mathāṇī}
  \item \emph{Say it:} \emph{mathāṇī}, once more
  \item \emph{Say it:} \emph{mathāṇī dahī mā̃y hai.} — The churning stick is in the curd.
  \item \emph{Recall:} say the Marwadi for to set curd, then the Marwadi for to churn, then say \emph{mathāṇī} again
\end{itemize}

\begin{tcolorbox}[breakable,colback=teal!4,colframe=teal!35!black,title={Before you move on}]
What does \mw{मथाणी} mean? (\textquotedblleft{}A churning stick\textquotedblright{}.) Say it once more.
\end{tcolorbox}
\section[\texorpdfstring{\mw{गाळणो} (gāḷṇo)}{गाळणो (gāḷṇo)}]{\mw{गाळणो} (gāḷṇo) — to strain}

\label{lesson:MW-C181-galno}



\begin{quote}
Before the new one: say the Marwadi for to churn, then the Marwadi for a churning stick.
\end{quote}

\subsection*{You'll want to know: \mw{गाळणो}}
\textbf{\mw{गाळणो}} — \emph{gāḷṇo} — \textquotedblleft{}to strain\textquotedblright{}.

\begin{quote}
\mw{मां} \mw{दूध} \mw{गाळै} \mw{है।}
\end{quote}

\begin{quote}
\emph{mā̃ dūdh gāḷai hai.}
\end{quote}

\begin{quote}
Mother strains the milk.
\end{quote}

Past (I …): \mw{म्हैं} \mw{पाणी} \mw{गाळ्यो}, \emph{mhaĩ pāṇī gāḷyo}.

\begin{cousinweb}
\textbf{\mw{गाळणो}} (\emph{gāḷṇo}) is inherited from Sanskrit \emph{gālayati}, \textquotedblleft{}filters, strains\textquotedblright{}; CDIAL lists Gujarati \emph{gāḷvũ} and Hindi \emph{gārnā} (CDIAL 4144). CDIAL gives no Marwari form in this entry.
\end{cousinweb}

\begin{grammarlens}[title={What you've built}]
One more word for this chapter.
\end{grammarlens}

\subsection*{Your turn}
\begin{itemize}
\raggedright
  \item \emph{Say it:} \emph{gāḷṇo}
  \item \emph{Say it:} \emph{gāḷṇo}, once more
  \item \emph{Say it:} \emph{mā̃ dūdh gāḷai hai.} — Mother strains the milk.
  \item \emph{Recall:} say the Marwadi for to churn, then the Marwadi for a churning stick, then say \emph{gāḷṇo} again
\end{itemize}

\begin{tcolorbox}[breakable,colback=teal!4,colframe=teal!35!black,title={Before you move on}]
What does \mw{गाळणो} mean? (\textquotedblleft{}To strain\textquotedblright{}.) Say it once more.
\end{tcolorbox}
\section[\texorpdfstring{\mw{खीर} (khīr)}{खीर (khīr)}]{\mw{खीर} (khīr) — rice pudding}

\label{lesson:MW-C181-khir}



\begin{quote}
Before the new one: say the Marwadi for a churning stick, then the Marwadi for to strain.
\end{quote}

\subsection*{You'll want to know: \mw{खीर}}
\textbf{\mw{खीर}} — \emph{khīr} — \textquotedblleft{}rice pudding\textquotedblright{}.

\begin{quote}
\mw{आज} \mw{मां} \mw{खीर} \mw{बणावै} \mw{है।}
\end{quote}

\begin{quote}
\emph{āj mā̃ khīr baṇāvai hai.}
\end{quote}

\begin{quote}
Mother makes rice pudding today.
\end{quote}

A feminine noun.

\begin{cousinweb}
\textbf{\mw{खीर}} (\emph{khīr}) is inherited from Sanskrit \emph{kṣīrī}, \textquotedblleft{}a dish made with milk\textquotedblright{}; CDIAL lists Hindi and Gujarati \emph{khīr}, rice boiled in milk (CDIAL 3705). CDIAL gives no Marwari form in this entry.
\end{cousinweb}

\begin{grammarlens}[title={What you've built}]
One more word for this chapter.
\end{grammarlens}

\subsection*{Your turn}
\begin{itemize}
\raggedright
  \item \emph{Say it:} \emph{khīr}
  \item \emph{Say it:} \emph{khīr}, once more
  \item \emph{Say it:} \emph{āj mā̃ khīr baṇāvai hai.} — Mother makes rice pudding today.
  \item \emph{Recall:} say the Marwadi for a churning stick, then the Marwadi for to strain, then say \emph{khīr} again
\end{itemize}

\begin{tcolorbox}[breakable,colback=teal!4,colframe=teal!35!black,title={Before you move on}]
What does \mw{खीर} mean? (\textquotedblleft{}Rice pudding\textquotedblright{}.) Say it once more.
\end{tcolorbox}
